\documentclass{aleph-notas}
\usepackage[utf8]{inputenc}
\usepackage[T1]{fontenc}
\usepackage{geometry}
\usepackage{graphicx}
\usepackage{float}
\usepackage{hyperref}
\usepackage[english,indonesian]{babel}

\title{Optimasi Pemrosesan Template LaTeX Menggunakan Pendekatan Berbasis AI}
\author{Academic Support Agent}
\date{\today}

\begin{document}
\maketitle

\begin{abstract}

\end{abstract}

\chapter{PENDAHULUAN}


\chapter{TINJAUAN PUSTAKA}


% Original content (commented out):
% % \documentclass[compacto,10pt]{aleph-notas}
\documentclass[compacto,10pt,comentarios]{aleph-notas}

% Se recomienda leer la documentación de esta
% clase en https://github.com/alephsub0/LaTeX_aleph-notas

% -- Paquetes adicionales
\usepackage{enumitem}
\usepackage{aleph-comandos}


% -- Datos del libro
\institucion{Proyecto Alephsub0}
% \carrera{}
\asignatura{Geometría Anaílica}
\tema{Modelo}
\autor{Andrés Merino}
\fecha{Abril 2024}
\fuente{montserrat}

%% --> Logos de las guias
\logouno[4.5cm]{Logos/LogoAlephsub0-02}
\definecolor{colordef}{cmyk}{0.81,0.62,0.00,0.22}

\begin{document}

\encabezado

\section{Ecuaciones trigonométricas}

\begin{advertencia}
    Suponemos conocidas las propiedades de la función 
    \[
        \func{\cos}{\R}{\R}
        \texty
        \func{\arccos}{[-1,1]}{\R}.
    \]
\end{advertencia}

Recuerde el siguiente resultado.

\begin{teo}
    Sean $x\in \R$ y $y\in [-1,1]$. Se tiene que si
    \[
        y =\cos(x),
    \]
    entonces
    \[
        x = \arccos(y) + 2k\pi
        \texto
        x = -\arccos(y) + 2k\pi
    \]
    con $k\in\Z$.
\end{teo}

%%%%%%%%%%%%%%%%%%%%%%%%%%%%%%%%%%%%%%%%%%%%%%%%%%%%
%% Ejercicio 1
%%%%%%%%%%%%%%%%%%%%%%%%%%%%%%%%%%%%%%%%%%%%%%%%%%%%
\begin{ejer}
    Resolver la ecuación
    \[
        1=4\cos\left(\frac{x}{3}\right).
    \]
\end{ejer}

\begin{proof}[Solución]
    Tomemos la ecuación y dividamos entre 4, obtenemos la expresión equivalente:
    \[
        \frac{1}{4}=\cos\left(\frac{x}{3}\right),
    \]
    con lo cual, las soluciones son
    \begin{enumerate}[label=\textit{\alph*)}]
        \item\label{ej01:c01} $\displaystyle \frac{x}{3}=\arccos\left(\frac{1}{4}\right)+2k\pi$, con $k\in\Z$; o
        \item\label{ej01:c02} $\displaystyle \frac{x}{3}=-\arccos\left(\frac{1}{4}\right)+2k\pi$, con $k\in\Z$.
    \end{enumerate}
    Así, la solución de la ecuación es
    \[
        x=3\arccos\left(\frac{1}{4}\right)+6k\pi
        \qquad\text{o}\qquad
        x=-3\arccos\left(\frac{1}{4}\right)+6k\pi.\qedhere
    \]
\end{proof}




\end{document}

\end{document}